%% abtex2-modelo-slides.tex, v-1.0 gfabinhomat
%% Copyright 2012-2015 by abnTeX2 group at http://www.abntex.net.br/ 
%%
%% This work may be distributed and/or modified under the
%% conditions of the LaTeX Project Public License, either version 1.3
%% of this license or (at your option) any later version.
%% The latest version of this license is in
%%   http://www.latex-project.org/lppl.txt
%% and version 1.3 or later is part of all distributions of LaTeX
%% version 2005/12/01 or later.
%%
%% This work has the LPPL maintenance status `maintained'.
%% 
%% The Current Maintainer of this work is Fábio Rodrigues Silva, 
%% member of abnTeX2 team, led by Lauro César Araujo. 
%% Further information are available on 
%% http://www.abntex.net.br/
%%
%% This work consists of the files abntex2-modelo-slides.tex, 
%% abntex2-modelo-references.bib and abntex2-modelo-marca.pdf
%%
%% Modelo desenvolvido por Fábio Rodrigues Silva (gfabinhomat@gmail.com)
%% Mais informações podem ser obtidas no guia do usuário Beamer 
%% (http://linorg.usp.br/CTAN/macros/latex/contrib/beamer/doc/beameruserguide.pdf)
%% Informações rápidas podem ser acessadas em http://en.wikibooks.org/wiki/LaTeX/Presentations


% Apresentações em widescreen. Outros valores possíveis: 1610, 149, 54, 43 e 32.
% Por padrão, as apresentações são no formato 4:3 (sem o aspectratio).
\documentclass[aspectratio=169]{beamer}	 	

\usetheme{Berlin}
\usecolortheme{dolphin}
\usefonttheme[onlymath]{serif}			% para fontes matemáticas
% Enconte mais temas e cores em http://www.hartwork.org/beamer-theme-matrix/ 
% Veja também http://deic.uab.es/~iblanes/beamer_gallery/index.html

% Customizações de Cores: fg significa cor do texto e bg é cor do fundo
% ---
% PACOTES
% ---
\usepackage[alf]{abntex2cite}		% Citações padrão ABNT
\usepackage[brazil]{babel}		% Idioma do documento
\usepackage{color}			% Controle das cores
\usepackage[T1]{fontenc}		% Selecao de codigos de fonte.
\usepackage{graphicx}			% Inclusão de gráficos
\usepackage[utf8]{inputenc}		% Codificacao do documento (conversão automática dos acentos)
\usepackage{txfonts}			% Fontes virtuais
% ---
\usepackage{caption}
% --- Informações do documento ---
\title{Demonstração do Controle por Modelo de Referência com Adaptação de Segundo Nível}
\author{Pedro Yochinori Gushiken}
\institute{Universidade Federal do Rio Grande do Norte}
\date{\today}
% ---

% ----------------- INÍCIO DO DOCUMENTO --------------------------------------
\begin{document}

% ----------------- NOVO SLIDE --------------------------------
\begin{frame}

\begin{minipage}{1\linewidth}
  \centering
\end{minipage}

\titlepage

\end{frame}

% ----------------- NOVO SLIDE --------------------------------
\begin{frame}{Sumário}
\tableofcontents
\end{frame}

% ----------------- NOVO SLIDE --------------------------------
\section{Introdução}
\begin{frame}
\frametitle{Contextualização}
\begin{enumerate}
\item<1-> Controle por Modelo de Referência\\
\item<2-> Identificação, Caso Indireto\\ 
\item<3-> Leis Adaptativas, Não-Linearidade\\
\item<4-> Calibragem dos Ganhos Adaptativos\\
\item<5-> Desempenho\\
\end{enumerate}
\end{frame}


% ----------------- NOVO SLIDE --------------------------------
\section{Conceitos}
\subsection{Problema do Controle Indireto Usando 1 Modelo de Identificação}

\begin{frame}
\frametitle{Descrição da Planta}

Seja a planta descrita pela equação de espaço de estado:
\begin{equation}
\dot{x}_p(t) = A_p x_p(t) + b u(t)
\end{equation}

Onde todas as variáveis de estado da planta são acessíveis, com $b=e_n$ e $A_p$ na seguinte forma canônica:

\[ 
A_p =
 \begin{bmatrix}
    0       & 1       &  0       &  ... &  0      \\
    0       & 0       &  1       &  ... &  0      \\
    0       & 0       &  0       &  ... &  0      \\
    ...     & ...     &  ...     &      &  ...    \\
    a_1     & a_2     &  a_3     &  ... &  a_n
  \end{bmatrix}
\]

Onde ${\theta_p}^*$ é o vetor transposto da última linha de $A_p$.  

\end{frame}

\begin{frame}
\frametitle{Descrição do Modelo de Referência}


Seja o modelo de referência descrito pela equação de espaço de estado:
\begin{equation}
\dot{x}_m(t) = A_m x_m(t) + b r(t)
\end{equation}

Com $b=e_n$ e $A_m$ na seguinte forma canônica:

\[ 
A_p =
 \begin{bmatrix}
    0       & 1       &  0       &  ... &  0      \\
    0       & 0       &  1       &  ... &  0      \\
    0       & 0       &  0       &  ... &  0      \\
    ...     & ...     &  ...     &      &  ...    \\
    {a_m}_1     & {a_m}_2     &  {a_m}_3     &  ... &  {a_m}_n
  \end{bmatrix}
\]

Onde o vetor $\theta_m$ transposto que corresponde à última linha da matriz $A_m$ é escolhido de forma que o modelo de referência seja estável.

\end{frame}

\begin{frame}
\frametitle{Declaração do Problema}

Queremos uma entrada $u(t)$ para a planta tal que a norma da diferença entre os vetores de variáveis de estado tenda a 0 na medida em que o tempo t tende a infinito. O método proposto de adaptação é indireto. No controle adaptativo indireto usando apenas um modelo, usamos as seguintes equações:

\begin{equation}
\dot{\hat{x}}_p = A_m\hat{x}_p(t) + [A_1(t)-A_m]x_p(t) +bu(t)
\end{equation}

É o modelo da planta, onde $A_1$ está na mesma forma canônica das matrizes mostradas anteriormente e sua última linha é o vetor $\theta_1$ transposto de estimativas dos parâmetros da planta, e $\phi_1(t)$ é o vetor de erros na estimativa de ${\theta_p}^*$.

\end{frame}

\begin{frame}
\frametitle{Equações, Cont.}

\begin{equation}
u(t) = r(t) +k^T(t)x_p(t)
\end{equation}

É a lei de controle.

\begin{equation}
k(t) = \theta_m - \theta_1 = k^* - \phi_1(t)
\end{equation}

É a parcela ajustada de forma adaptativa do controlador. Onde $k^*$ é o vetor de parâmetros do controlador caso a planta fosse conhecida e $\phi_1(t)$ é o vetor de erros entre os parâmetros da planta e as estimativas dos parâmetros.

\end{frame}
\begin{frame}
\frametitle{Lei Adaptativa}

Definindo o erro de identificação $e_1$ como:
\begin{equation}
e_1(t) = \hat{x}_p(t) - x_p(t)
\end{equation}

Temos a seguinte lei adaptativa:

\begin{equation}
\dot{\theta_1} = \dot{\phi_1} = -{e_1}^TPbx_p(t)
\end{equation}

A partir da teoria de Lyapunov usando a candidata $V(e_1(t), \phi_1(t))= \frac{1}{2} {e_1}^T(t)Pe_1(t) + {\phi_1}^T\phi_1(t)$ é possível provar que o erro de controle $e_1$ tende a 0 para o tempo t tendendo a infinito para qualquer referência e que o erro de estimativa dos parâmetros $\phi_1$ tende a 0 para t tendendo a infinito para r(t) que é persistentemente excitante.


\end{frame}

\subsection{Controle Indireto com Adaptação de Segundo Nível (Múltiplos Modelos de Identificação)}

\begin{frame}
\frametitle{Localização de ${\theta_p}^*$}

Em \cite{han1}, foi provado que, se o vetor de parâmetros desconhecidos ${\theta_p}^*$ da planta a ser controlada estiver contido no casco convexo dos parâmetros iniciais $\theta_i(0)$ de N modelos (onde $N\geq n+1$), dentro de certas condições ${\theta_p}^*$ continuará contido dentro do casco convexo delimitado pelos vetores $\theta_i(t)$ de cada modelo a medida que estes de adaptam.

\end{frame}

\begin{frame}
\frametitle{Perguntas Decorrentes}

Assumindo que o vetor de parâmetros desconhecidos da planta ${\theta_p}^*$ se encontra em alguma região conhecida, as seguintes questões se seguem:

a)Quantos Modelos são necessários para identificar ${\theta_p}^*$?\\
b)Estes Modelos devem ser fixos ou adaptativos?\\
c)Em ambos os casos, como seriam escolhidos os parâmetros iniciais de cada modelo?\\

\end{frame}

\begin{frame}
\frametitle{Regiões Convexas/ Cascos Convexos}
Para delimitar uma região convexa do $\Re^n$ cuja dimensão seja $\Re^n$ são necessários no mínimo $n+1$ pontos naquele espaço.\\

Qualquer ponto no interior de uma região convexa pode ser expresso como uma combinação linear dos vetores que delimitam a região cumprindo as seguintes condições:\\ 

a) A soma dos coeficientes da combinação seja igual a 1\\ 
b) Todos os coeficientes maiores que 0\\

\end{frame}

\begin{frame}
\frametitle{Exemplo 1: Região Convexa}

\centering
\begin{figure}
  \includegraphics[width=.69\linewidth]{./pics/Ex_Regiao.png}
    \captionof{figure}{Em roxo uma região convexa do $\Re^2$, delimitada por [2 5], [6 1] e [9 4]}
    \label{fig:test2912}
\end{figure}


\end{frame}
\begin{frame}
\frametitle{Múltiplos Modelos}

Tanto no caso de modelos fixos quanto adaptativos, a equação diferencial que descreve o modelo é a mesma. Com $x_i(t_0) = x_p(t_0)$. $A_i$ se encontra na mesma forma canônica descrita anteriormente. Desta forma o modelo $\Sigma_i$ é descrito por:

\begin{equation}
\Sigma_i: \dot{\hat{x}}_i = A_m\hat{x}_p(t) + [A_i(t)-A_m]x_p(t) +bu(t)
\end{equation}

No caso em que os modelos são adaptativos, os vetores de parâmetros $\theta_i$ são atualizados usando as leis adaptativas descritas anteriormente, com $e_i(t) =x_i(t)-x_p(t) $. No caso em que os modelos são fixos, estes vetores são constantes, e portanto os modelos não são estritamente de identificação, porém mesmo assim ainda é possível estimar ${\theta_p}^*$.

\end{frame}

\begin{frame}

Tendo em mente para o caso adaptativo que:


\centering
\begin{minipage}{.48\textwidth}
  \centering
\begin{equation}
\sum_{i=1}^{n+1}\alpha_i\theta_i(t_0) = \theta_p
\end{equation}
\begin{equation}
\dot{e}_i = A_m e_i(t) + b{\phi_i}^T(t)x_p(t)
\end{equation}
\begin{equation}
e_i(t_0)=0
\end{equation}
\end{minipage}%
  \hfill
\begin{minipage}{.48\textwidth}
  \centering

\begin{equation}
\dot{\theta}_i(t) = \dot{\phi}_i(t) = -{e_i}^T(t)Pbx_p(t)
\end{equation}
\begin{equation}
\theta_i(t_0) = \theta_{i0}
\end{equation}
\begin{equation}
\sum_{i=1}^{n+1}\alpha_i\phi(t_0)=0
\end{equation}
\end{minipage}


\end{frame}


\begin{frame}

É possível enxergar as constantes $\alpha_1,\alpha_2,...,\alpha_{n+1}$ como uma alternativa de parametrização da planta. Consequentemente, a velocidade e precisão na estimação de $\alpha_i$ é que determinam seu desempenho em comparação com a adaptação de primeiro nível. Dadas as equações anteriores se seguem as seguintes equações:
\centering
\begin{minipage}{.48\textwidth}
  \centering
  \begin{equation}
  \sum_{i=1}^{n+1}\alpha_i\theta_i(t) = \theta_p \space 
  \end{equation}
\end{minipage}%
  \hfill

\begin{minipage}{.48\textwidth}
  \centering
    \begin{equation}
    \sum_{i=1}^{n+1}\alpha_i e_i(t) = 0 \space
    \end{equation}
\end{minipage}

\end{frame}

\begin{frame}
E em forma matricial:

\begin{equation}
\bar{\Theta}(t)\bar{\alpha} = \theta_p
\end{equation}
\begin{equation}
\bar{E}(t)\bar{\alpha} = 0
\end{equation}

Onde $\bar{\Theta}$ e $\bar{E}$ são matrizes $n\times(n+1)$.

\end{frame}

\begin{frame}
E devido às propriedades da região convexa, podemos expressar:
\begin{equation}
\Theta(t)\alpha = \theta_p - \theta_{n+1}
\end{equation}
\begin{equation}
E(t)\alpha = -e_{n+1}(t)
\end{equation}

Onde $\Theta(t)$ e $E(t)$ são matrizes $n\times n$, onde cada vetor coluna é $\theta_i(t)-\theta_{n+1}(t)$ e $e_i(t)-e_{n+1}(t)$ respectivamente e $\alpha^T = [\alpha_1,...,\alpha_n]$. Finalmente, substituindo $\sum_{i=1}^{n+1}\alpha_i\theta_p = \theta_p$ em (19), obtemos:

\begin{equation}
\Phi(t)\alpha = \phi_{n+1}(t)
\end{equation}

Onde $\Phi(t)$ é uma matriz $n\times n$ cujas colunas são os vetores de erros de parâmetro até $i=n$.

\end{frame}

\subsection{Computação de $\alpha$}
\begin{frame}
\frametitle{Computação Algébrica usando (19)}
Como $\theta_p$ e $\alpha$ são ambos desconhecidos, os vetores $\theta_i$ devem ser avaliados em instantes de tempo distintos, para estimar $\alpha$. Se as medidas são feitas nos tempos $t_1$ e $t_2$ obtemos:

\begin{equation}
[\Theta(t_1)-\Theta(t_2)]\alpha =-[\theta_{n+1}(t_1) - \theta_{n+1}(t_2)]
\end{equation}

A partir da qual podemos computar $\alpha$. Alternativamente, fazendo uma operação de diferenciação sobre os dois lados de (19) obtemos:

\begin{equation}
\dot{\Theta}(t)\alpha = -\dot{\theta}_{n+1}(t)
\end{equation}

\end{frame}

\begin{frame}
Como $\dot{\theta_i}$ são as leis adaptativas, chegamos novamente a:
\begin{equation}
E(t)\alpha = -e_{n+1}(t)
\end{equation}

Como $E(t)$ é obrigatoriamente uma matriz invertível, $\alpha$ também pode ser computado a partir de (24) usando a inversa de $E(t)$.  

\end{frame}

\begin{frame}
\frametitle{Comentário sobre Múltiplos Modelos Fixos}

Embora para um único modelo fixo não seja possível identificar o vetor de parâmetros $\theta_p$ da planta, como as equações de erro de identificação são lineares obtemos (24) novamente, relacionando os erros de saída $e_i(t)$ à $\alpha$. Desta forma, coletivamente, modelos fixos podem ser utilizados para identificar uma planta desde que tenhamos algum conhecimento prévio da região em que estão localizados seus parâmetros.
 
\end{frame}


% ----------------- NOVO SLIDE --------------------------------
\section{Simulações com Modelos Fixos}
% % % % % % % % % % % % % % % % % % % % % % % % % % % % % % % % % % %
\subsection{Simulação 1}
\begin{frame}
\frametitle{Dados da Simulação}

\begin{equation}
\theta_p^T = [-2~~~-1]
\end{equation}

\begin{equation}
\theta_m^T = [-1~~~-3]
\end{equation}

\begin{equation}
\theta_1^T = [-2~~~-4]
\end{equation}

\begin{equation}
\theta_2^T = [-6~~~-6]
\end{equation}

\begin{equation}
\theta_3^T = [3~~~-6]
\end{equation}

\end{frame}

\begin{frame}
\frametitle{Erro de Saída, Norma do Erro de Saída}
\centering
\begin{minipage}{.48\textwidth}
  \centering
  \includegraphics[width=\linewidth]{./pics/ES_SIM1.png}
  \captionof{figure}{Erro de Saída}
\end{minipage}%
  \hfill
\begin{minipage}{.48\textwidth}
  \centering
  \includegraphics[width=\linewidth]{./pics/Norm2_ES_SIM1.png}
  \captionof{figure}{Norma 2 do Erro de Saída}
\end{minipage}
\end{frame}

\begin{frame}
\frametitle{Saída}
  \centering
  \includegraphics[width=0.7\linewidth]{./pics/S_MeP_SIM1.png}
  \captionof{figure}{Saída da Planta e do Modelo}
\end{frame}

\subsection{Simulação 2}

\begin{frame}
\frametitle{Dados da Simulação}

\begin{equation}
\theta_p^T = [2~~~1]
\end{equation}

\begin{equation}
\theta_m^T = [-1~~~-3]
\end{equation}

\begin{equation}
\theta_1^T = [3~~~-6]
\end{equation}

\begin{equation}
\theta_2^T = [-2~~~-2]
\end{equation}

\begin{equation}
\theta_3^T = [7~~~-3]
\end{equation}

\end{frame}


\begin{frame}
\frametitle{Erro de Saída, Norma do Erro de Saída}
\centering
\begin{minipage}{.48\textwidth}
  \centering
  \includegraphics[width=\linewidth]{./pics/ES_SIM2.png}
  \captionof{figure}{Erro de Saída}
\end{minipage}%
  \hfill
\begin{minipage}{.48\textwidth}
  \centering
  \includegraphics[width=\linewidth]{./pics/Norm2_ES_SIM2.png}
  \captionof{figure}{Norma 2 do Erro de Saída}
\end{minipage}

\end{frame}

\begin{frame}
\frametitle{Saída}
  \centering
  \includegraphics[width=0.7\linewidth]{./pics/S_MeP_SIM2.png}
  \captionof{figure}{Saída da Planta e do Modelo}
\end{frame}

\subsection{Simulação 3}

\begin{frame}
\frametitle{Dados da Simulação}

\begin{equation}
\theta_p^T = [5~~~3]
\end{equation}

\begin{equation}
\theta_m^T = [-24~~~-8]
\end{equation}

\begin{equation}
\theta_1^T = [-10~~~-10]
\end{equation}

\begin{equation}
\theta_2^T = [15~~~-10]
\end{equation}

\begin{equation}
\theta_3^T = [5~~~15]
\end{equation}

\end{frame}


\begin{frame}
\frametitle{Erro de Saída, Norma do Erro de Saída}
\centering
\begin{minipage}{.48\textwidth}
  \centering
  \includegraphics[width=\linewidth]{./pics/ES_SIM3.png}
  \captionof{figure}{Erro de Saída}
\end{minipage}%
  \hfill
\begin{minipage}{.48\textwidth}
  \centering
  \includegraphics[width=\linewidth]{./pics/Norm2_ES_SIM3.png}
  \captionof{figure}{Norma 2 do Erro de Saída}
\end{minipage}
\end{frame}

\begin{frame}
\frametitle{Saída}
  \centering
  \includegraphics[width=0.7\linewidth]{./pics/S_MeP_SIM3.png}
  \captionof{figure}{Saída da Planta e do Modelo}
\end{frame}

\subsection{Simulação 4}

\begin{frame}
\frametitle{Dados da Simulação}

\begin{equation}
\theta_p^T = [3+sin(t)~~~4+cos(t)~~~3*]
\end{equation}

\begin{equation}
\theta_m^T = [-15~~~-23~~~-9]
\end{equation}

\begin{equation}
\theta_1^T = [5~~~3~~~3]
\end{equation}

\begin{equation}
\theta_2^T = [1~~~3~~~3]
\end{equation}

\begin{equation}
\theta_3^T = [3~~~6~~~3]
\end{equation}

\end{frame}


\begin{frame}
\frametitle{Erro de Saída, Norma do Erro de Saída}
\centering
\begin{minipage}{.48\textwidth}
  \centering
  \includegraphics[width=\linewidth]{./pics/ES_SIM4.png}
  \captionof{figure}{Erro de Saída}
\end{minipage}%
  \hfill
\begin{minipage}{.48\textwidth}
  \centering
  \includegraphics[width=\linewidth]{./pics/Norm2_ES_SIM4.png}
  \captionof{figure}{Norma 2 do Erro de Saída}
\end{minipage}
\end{frame}

\begin{frame}
\frametitle{Saída}
  \centering
  \includegraphics[width=0.7\linewidth]{./pics/S_MeP_SIM4.png}
  \captionof{figure}{Saída da Planta e do Modelo}
\end{frame}

% ----------------- NOVO SLIDE --------------------------------
\section{Referências}

% --- O comando \allowframebreaks ---
% Se o conteúdo não se encaixa em um quadro, a opção allowframebreaks instrui 
% beamer para quebrá-lo automaticamente entre dois ou mais quadros,
% mantendo o frametitle do primeiro quadro (dado como argumento) e acrescentando 
% um número romano ou algo parecido na continuação.
% --- O comando \allowframebreaks ---
% Se o conteúdo não se encaixa em um quadro, a opção allowframebreaks instrui 
% beamer para quebrá-lo automaticamente entre dois ou mais quadros,
% mantendo o frametitle do primeiro quadro (dado como argumento) e acrescentando 
% um número romano ou algo parecido na continuação.

\begin{frame}[allowframebreaks]
\frametitle{Referências}
\nocite{narendra1}
\nocite{narendra_main}
\nocite{narendra2012rationale}
\nocite{narendra1997}
\nocite{han1}
\nocite{han2010}
\nocite{han2012}
\nocite{kuipers1}
\nocite{ioannou2}
\nocite{chen}
\nocite{ogata}
\nocite{khalil}
\end{frame}
% ----------------- FIM DO DOCUMENTO -----------------------------------------
\bibliography{bib-controleadaptativo}
\end{document}
